% ============================================================================
% Configuración y comandos personalizados de la memoria
% ============================================================================

% -----------------------------
% Datos de la memoria
% -----------------------------
\newcommand{\TituloMemoria}{Transformers para la predicción de series de tiempo no equiespaciadas}
\newcommand{\AutorNombre}{Gerald Méndez}
\newcommand{\GradoQueOpta}{Ingeniero Civil Matemático}
\newcommand{\Carrera}{Ingeniería Civil Matemática}
\newcommand{\ProfesorGuia}{Carlos Valle}
\newcommand{\ProfesorCoguia}{Francisco Cuevas} % deja en blanco o borra si no tienes
\newcommand{\Institucion}{Universidad Técnica Federico Santa María}
% \newcommand{\Facultad}{Departamento de }
\newcommand{\Departamento}{Departamento de Matemáticas}
\newcommand{\Ciudad}{Valparaíso}
\newcommand{\Anio}{2025}

% Si quieres usar el logo de la universidad en la portada:
% Coloca el archivo en la carpeta images/ y ajusta el nombre aquí.
\newcommand{\LogoUniversidad}{logo_universidad} % sin extensión; se asume .pdf, .png, etc.

% -----------------------------
% Comandos matemáticos útiles
% -----------------------------
\newcommand{\R}{\mathbb{R}}
\newcommand{\N}{\mathbb{N}}
\newcommand{\Z}{\mathbb{Z}}
\newcommand{\E}{\mathbb{E}}
\newcommand{\Var}{\mathrm{Var}}
\newcommand{\Cov}{\mathrm{Cov}}

\newcommand{\vect}[1]{\mathbf{#1}}
\newcommand{\mat}[1]{\mathbf{#1}}
\newcommand{\ts}[1]{\boldsymbol{#1}} % para series de tiempo

% Notación específica del trabajo (puedes ir agregando más)
\newcommand{\Transformer}{\textit{Transformer}}
\newcommand{\Transformers}{\textit{Transformers}}
\newcommand{\CTTransformer}{CT-Transformer}
\newcommand{\CTTransformerFull}{Continuous-Time Transformer (\CTTransformer)}
\newcommand{\CTTransformerBase}{\CTTransformer\ base}
\newcommand{\CTTransformerSmall}{\CTTransformer\ pequeño}
\newcommand{\CTTransformerLarge}{\CTTransformer\ grande}

% -----------------------------
% Estilos de teoremas (si los necesitas)
% -----------------------------
\theoremstyle{plain}
\newtheorem{teorema}{Teorema}[chapter]
\newtheorem{proposicion}[teorema]{Proposición}
\newtheorem{corolario}[teorema]{Corolario}

\theoremstyle{definition}
\newtheorem{definicion}[teorema]{Definición}
\newtheorem{ejemplo}[teorema]{Ejemplo}

\theoremstyle{remark}
\newtheorem{observacion}[teorema]{Observación}

% -----------------------------
% Otros ajustes de estilo opcionales
% -----------------------------
% Espaciado entre líneas, si tu escuela lo permite:
% \usepackage{setspace}
% \onehalfspacing
